% =====================================================================
% §4 A STRUCTURAL INTERPRETATION — consolidation of the former
% Model / Equilibrium Analysis / Identification / Results sections.
% 4.1 is new prose; 4.2-4.5 are relocated verbatim, demoted one level.
% Labels sec:model, sec:equilibrium, sec:identification, sec:results are
% carried onto the corresponding subsections (they have live references).
% =====================================================================

\section{A Structural Interpretation}
\label{sec:model}

\subsection{Motivation}
\label{ss:model_motivation}

To recover the forces behind the patterns documented in
Section~\ref{sec:setting}, we specify and estimate a model of the seller's disclosure
decision. The model is a housing-market application of \citet{dye_disclosure_1985}, as
equilibrium-characterised by \citet{jung_kwon_1988}, with one structural addition:
heterogeneity in buyer EPC sensitivity ($\rho$). The binary quality space
($\theta\in\{0,1\}$) and the seller-information heterogeneity ($\pi$) are adapted directly
from those antecedents. What goes beyond them is the closed-form expression for the
non-disclosure premium and its invariance to market composition in the partial-pooling
regime (Proposition~\ref{prop:equilibrium}), the strict cross-regime ordering of the premium
(Eq.~\ref{eq:premium}), and an identification strategy that uses the directly observed audit
status $K$ to recover $c_0$ and $\rho$ market by market.

We allow for, but do not impose, strategic concealment. Nothing in the specification requires
informed sellers to withhold: the hiding rate is an equilibrium object, the model admits full
unravelling as a special case, and whether concealment occurs --- and how much --- is what the
estimation determines.

Two considerations motivate this structure. The first is resemblance. The facts in
Section~\ref{sec:facts} line up with the Dye--Jung-Kwon configuration: a silent pool that is
overwhelmingly silent by constraint rather than by choice, a small but non-trivial mass of
informed sellers who withhold, and a premium attached to silence rather than a penalty. The
second is institutional. The binary quality space has an empirical anchor in the rental-ban
threshold introduced by \emph{loi Climat et R\'esilience}: dwellings rated E, F, or G face a
progressive prohibition on new lettings, which draws a bright line between the bottom of the
grade distribution and the rest and maps onto $\theta\in\{0,1\}$ without an arbitrary cutoff.

A competing explanation must also clear a demanding timing test. Any account that attributes
the post-reform widening to a supply-side shift has to generate an abrupt change at
July~2021 in \emph{both} the quality premium and the non-disclosure premium, in the same
markets and on the same calendar date, while leaving the composition of the listing pool
undisturbed. The share of listings without a disclosed score is, however, essentially flat
across the reform in the apartment segment (Section~\ref{ss:fact2}), which is the segment
where identification is sharpest. A shift in the mix of sellers would generically move the
two gaps by different amounts, and would show up in that share. Neither is observed.

The mechanism itself is the one set out in Section~\ref{sec:unraveling}: sellers who cannot
disclose dilute the adverse inference buyers draw from silence, which makes withholding
attractive to informed sellers of poor quality, whose entry the buyer cannot separate from
the involuntarily silent. France sits in the bottom-right cell of
Figure~\ref{fig:equilibrium_outcomes} --- incomplete registry coverage combined with positive
buyer EPC-sensitivity. Even where the EPC-sensitive buyer share approaches one, as it does in
most apartment markets after the reform, incomplete coverage alone sustains the premium: so
long as a non-trivial share of non-disclosers hold no certificate, buyers cannot fully
penalise silence. Full unravelling requires both conditions; France satisfies only one.

Finally, we are explicit about what the model does not settle. It takes no stance on
\emph{why} an informed seller withholds --- deliberate concealment and an unrecorded but
still-valid certificate are observationally similar in the listing data, and we do not
attempt to separate them. It assumes no falsification, so a displayed score is taken at face
value. And it maintains Corollary~\ref{cor:hq_disclose}, $\sigma_1^*=0$, rather than testing
it: informed non-disclosers are assigned low quality by revealed preference, which is an
assumption and not a result. Appendix~\ref{app:robustness} reports estimates that instead read
the withheld quality from the registry, so the restriction can be assessed. Nor does the
design pin down $\rho_m$ from the interior in most markets: the data are consistent with full
buyer sensitivity but cannot establish how close to one it lies, and we maintain
$\hat\rho_m = 1$ for those markets and treat $\hat c_{0,m}$ as an upper bound accordingly. We
do not purport to estimate the one true model of housing disclosure; we propose one way of
reading these data through a workhorse framework.

\subsection{Setup}
\label{ss:model_setup}

\subsubsection{Environment and Timing}

A seller owns a property with observable attributes $X$ and an unobservable energy performance level $\theta \in \{0,1\}$, where $\theta=1$ denotes high efficiency. Whether the seller knows $\theta$ depends on whether she has obtained a certified Diagnostic de Performance \'{E}nerg\'{e}tique (DPE) audit prior to listing. With probability $\pi\in(0,1)$, the seller is \emph{informed} ($K=1$) and knows $\theta$; with probability $1-\pi$, she is \emph{uninformed} ($K=0$) and does not. The parameter $\pi$ is not a choice variable: it reflects structural features of local audit markets, including the cost and availability of accredited assessors, the time pressure sellers face (relocation, financial distress, inheritance), and the regulatory requirements in force. It is measured directly from the ADEME registry match, as described in Section~\ref{sec:data}.

The seller simultaneously chooses a disclosure decision $d\in\{0,1\}$ and a listing price $p^{\text{list}}$. French property law requires that all mandatory diagnostic documents, including the DPE, be appended to the sales agreement at signing (loi ALUR, Article L271-4 CCH), so the disclosure decision and the listing price are set jointly at the time the property is placed on the market. This institutional timing motivates the simultaneous-choice formulation: the seller does not first decide whether to disclose and then price, but sets a complete strategy that internalizes how the disclosure decision shifts buyer beliefs and thereby affects both the probability of sale and the expected transaction price.

\noindent\textbf{Assumption (No Falsification).} \textit{Disclosed EPC scores are verifiable against the ADEME registry. An informed seller with $K=1$ may choose $d\in\{0,1\}$, while an uninformed seller with $K=0$ is constrained to $d=0$: she cannot disclose what she does not possess. This rules out costless misreporting.}

\medskip
The key friction is that buyers observe $(d, p^{\text{list}}, X)$ but not the seller's information status $K$: they cannot distinguish a seller who chose not to disclose from one who could not. A buyer then arrives and decides whether to purchase at $p^{\text{list}}$. With probability $\rho\in(0,1)$, the buyer is \emph{EPC-sensitive} and incorporates energy performance into her valuation; with probability $1-\rho$, the buyer is \emph{EPC-insensitive} and conditions her bid solely on observable attributes. Sellers cannot observe buyer type ex ante and optimize against expected demand. The parameter $\rho$ is market-specific, reflecting the local composition of the buyer pool, and is identified from the differential price response to the July 2021 reform across markets, as detailed in Section~\ref{sec:identification}.

We provide a visual representation of the model dynamics and timing in Figure~\ref{fig:timing}; a fuller schematic overview appears later in Figure~\ref{fig:vis_model}.

\begin{figure}[H]
    \centering
    \includegraphics[width=\textwidth]{fig/visuals/model_dynamics.png}
    \caption{Model timing}
\label{fig:timing}
\end{figure}

\subsubsection{Payoffs and Hiding Costs}

In French residential markets, listing prices serve as anchors for negotiation rather than posted take-it-or-leave-it offers: the transaction price conditional on a sale typically falls below the listing price by a margin that depends on local market conditions and property characteristics.\footnote{The negotiation margin in France is empirically small: the majority of properties transact at or very near list price, and the average discount rarely exceeds three to five percent. \citet{merlo_bargaining_2004} document similar patterns for England, showing that most residential transactions involve little or no negotiation once a listing price is set. \citet{genesove_loss_2001} provide evidence that list prices act as anchors that constrain final transaction prices. For France specifically, \citet{friggit_marche_2013} documents that the average negotiation margin has remained narrow and declining, with a large share of properties selling without any price reduction.} Buyer arrival and match quality are stochastic, so the seller faces uncertainty over both whether a sale occurs and, conditional on a sale, at what price. The seller's problem therefore separates into a transaction probability $\Gamma$ and an expected transaction price $\phi$, each depending on the listing price and disclosure decision:
\begin{equation}
    \max_{\{d,\,p^{\text{list}}\}} \;\Pi = \underbrace{\mathbb{E}\!\left[p^{\text{trans}}\mid d,p^{\text{list}},X\right]}_{\phi(p^{\text{list}})} \cdot \underbrace{\Pr(T{=}1\mid d,p^{\text{list}},X)}_{\Gamma(p^{\text{list}})} - C(\theta,d)
    \label{eq:payoff}
\end{equation}
where $\phi(p^{\text{list}}) = \gamma_0+\gamma_1 \cdot p^{\text{list}}+\sum_n\gamma_n X_n$ is the expected transaction price conditional on a sale occurring, and $C(\theta,d)=\mathbf{1}[d=0]\cdot(c_0+\theta c_1)$ is the hiding cost, paid only when the seller withholds information.

The hiding cost has two components. The term $c_0\geq 0$ applies to all non-disclosers, regardless of quality. In the French institutional context, $c_0$ captures the regulatory burden imposed by loi ALUR (legal exposure if a buyer subsequently discovers a concealed audit via the public ADEME registry, notaire liability checks, and the informational ambiguity that prolongs negotiation and increases sale failure risk). The term $c_1\geq 0$ is an additional opportunity cost borne exclusively by high-quality sellers ($\theta=1$) who choose to hide: by not disclosing, they forgo the revenue premium that verified high quality would command from EPC-sensitive buyers. This asymmetry, $c_1>0$, is economically decisive: it means that hiding is strictly cheaper for low-quality sellers than for high-quality ones, so the net gain from concealment is larger for $\theta=0$ sellers at any given pool belief. This differential is the sorting pressure that drives full separation of high-quality sellers in equilibrium (Corollary~\ref{cor:hq_disclose}). Hiding costs are paid unconditionally on sale: they represent compliance burdens and reputational exposure that arise at listing time, not at closing. A note on the dual role of $c_0$ is in order. As a structural primitive, $c_0$ is an institutional cost of concealment that we treat as fixed within the prevailing enforcement regime and recover from the equilibrium price gap (Section~\ref{sec:identification}); as a policy object, it is the margin a regulator shifts through stricter enforcement, and we study counterfactual variations in $c_0$ as such in Section~\ref{sec:policy}. The two uses are consistent: estimation reads $c_0$ off the current regime, while the policy exercise perturbs it around that baseline.

\subsubsection{Buyer Valuations and Demand}

The partition of buyers into EPC-sensitive and EPC-insensitive types reflects a well-documented empirical regularity: willingness to pay for energy efficiency varies sharply with how much the buyer will bear the ongoing energy costs. Primary-residence buyers in northern France with long heating seasons, and investor-buyers who face the impending F/G rental ban under the \textit{loi Climat et R\'{e}silience}, have strong incentives to incorporate EPC scores into their bids. Vacation-home buyers who occupy the property a few weeks per year, and buyers in mild climates, may rationally discount energy performance. This heterogeneity determines $\rho$, the share of the local buyer pool that is EPC-sensitive.

EPC-sensitive buyers ($B=1$) condition their valuation on energy performance:
\begin{equation}
    v_1(\theta,d) = \begin{cases}
        \beta_0 + \beta_1\cdot\theta + \textstyle\sum_n\beta_n X_n & \text{if } d=1 \;\text{(verified quality)},\\
        \beta_0 + \beta_1\cdot\mathbb{E}[\theta\mid d{=}0] + \textstyle\sum_n\beta_n \cdot X_n & \text{if } d=0 \;\text{(inferred quality)}.
    \end{cases}
\end{equation}
We adopt the benchmark $\alpha_k=\beta_k$: buyers value quality equally whether it is certified or inferred. This is a conservative assumption in a precise sense: under the alternative $\alpha_k<\beta_k$ (a distrust discount on unverified quality), the non-disclosure premium would be smaller for any given pool composition, so our identified $\rho$ is a lower bound on the true EPC-sensitive buyer share. The non-disclosure premium we document therefore understates what a fully skeptical buyer would generate. EPC-insensitive buyers ($B=0$) value the property solely on observables: $v_0(X)=\delta_0+\sum_n\delta_n X_n$. The aggregate transaction probability mixes the two buyer types:
\begin{equation}
    \Gamma(p^{\text{list}}) = \rho\cdot G\!\left(v_1(d)-p^{\text{list}}\right) + (1-\rho)\cdot G\!\left(v_0-p^{\text{list}}\right)
    \label{eq:demand}
\end{equation}
where $G(\cdot)$ is the logistic CDF governing idiosyncratic taste shocks. The parameter $\rho$ enters the demand function as a weight on the EPC-sensitive buyer valuation: a higher $\rho$ amplifies the revenue difference between a disclosed high-quality listing and a non-disclosing listing, raising the stakes of the disclosure decision for the seller. It is this amplification channel through which $\rho$ governs disclosure incentives. Empirically, $\rho$ is a market-level structural parameter identified from the cross-market variation in the price response to the July 2021 reform, as developed in Section~\ref{sec:identification}.

\subsubsection{Buyer Beliefs and the Non-Disclosing Pool}

When observing a non-disclosing listing, buyers cannot determine whether the seller is strategically concealing an unfavorable audit or simply lacks a certified EPC. They form a posterior belief about expected quality using Bayes' rule. The structure of this belief depends on the composition of the non-disclosing pool, which is itself an equilibrium object.

\paragraph{Who Hides?}

Each seller is characterized by two attributes: information status $K \in \{0,1\}$ and quality $\theta \in \{0,1\}$. We maintain the following assumption throughout.

\begin{assumption}[Independence]\label{ass:independence}
$K \perp \theta$, with $\Pr(K=1) = \pi$ and $\Pr(\theta=1) = \bar{\theta}$.
\end{assumption}

\noindent Independence is most plausible for transactions driven by exogenous life events (succession following a death, court-ordered liquidation, or divorce settlement), in which legal proceedings impose binding timing constraints that preclude any quality-contingent decision to obtain a certified EPC. In such cases, the seller's information status $K$ is determined by administrative capacity and procedural deadlines rather than the energy performance of the dwelling; a seller under a judicial mandate to liquidate an estate has neither the incentive nor the discretion to condition audit acquisition on $\theta$. In France, forced transactions of this type constitute a substantial share of the resale market, providing a natural mechanism through which $K \perp \theta$ obtains in the data.

Table~\ref{tab:who_hides} summarizes disclosure behavior by seller type. Uninformed sellers have no choice: they cannot disclose information they do not possess (No Falsification Assumption). Informed sellers choose strategically, weighing the revenue gain from disclosure against the option of pooling with the non-disclosing mass. We denote by $\sigma_\theta \equiv \Pr(d=0 \mid K=1, \theta)$ the equilibrium hiding probability for an informed seller with quality $\theta$.

\begin{table}[H]
\centering
\caption{Disclosure Behavior by Seller Type}
\label{tab:who_hides}
\small
\renewcommand{\arraystretch}{1.3}
\setlength{\tabcolsep}{20pt}
\begin{tabular}{lcccl}
\toprule
\toprule
\textbf{Type} & \textbf{Mass} & \textbf{Can Disclose?} & \textbf{Hides?} & \textbf{Hiding Mass} \\
\midrule
Informed, $\theta=1$ & $\pi \bar{\theta}$ & Yes & Prob.\ $\sigma_1$ & $\pi \bar{\theta} \sigma_1$ \\
Informed, $\theta=0$ & $\pi (1-\bar{\theta})$ & Yes & Prob.\ $\sigma_0$ & $\pi (1-\bar{\theta}) \sigma_0$ \\
Uninformed, $\theta=1$ & $(1-\pi) \bar{\theta}$ & No & Always & $(1-\pi) \bar{\theta}$ \\
Uninformed, $\theta=0$ & $(1-\pi) (1-\bar{\theta})$ & No & Always & $(1-\pi) (1-\bar{\theta})$ \\
\bottomrule
\bottomrule
\end{tabular}

\vspace{0.5em}
\raggedright
\footnotesize
\textit{Notes:} $\sigma_\theta \equiv \Pr(d=0 \mid K=1, \theta)$ denotes the hiding probability for an informed seller with quality $\theta$. Uninformed sellers cannot disclose and hide with probability one.
\end{table}

\paragraph{The Belief Function}

Applying Bayes' rule to the pool composition in Table~\ref{tab:who_hides}, the expected quality among non-disclosers is:
\begin{equation}
    e(\sigma_0, \sigma_1) \equiv \Pr(\theta=1 \mid d=0) = \frac{\pi \bar{\theta} \sigma_1 + (1-\pi) \bar{\theta}}{\pi \bar{\theta} \sigma_1 + \pi (1-\bar{\theta}) \sigma_0 + (1-\pi)}
    \label{eq:belief_general}
\end{equation}

The numerator counts the high-quality mass in the pool: informed high-quality sellers who choose to hide ($\pi\bar{\theta}\sigma_1$) together with uninformed sellers who happen to have high quality ($(1-\pi)\bar{\theta}$). The denominator is the total non-disclosing mass. The belief $e$ is therefore a weighted average across two economically distinct populations. When uninformed sellers dominate the pool ($\pi$ small or $\sigma_0$ small), non-disclosure looks like a constraint rather than a choice and beliefs remain favorable to the seller. When strategic low-quality hiders dominate ($\pi$ large and $\sigma_0$ large), the pool is contaminated and beliefs deteriorate. This is why unraveling is incomplete: non-disclosure is inherently ambiguous when $\pi < 1$, whereas in classical disclosure theory all sellers are informed and silence is unambiguous bad news.

\begin{lemma}[Belief Monotonicity]\label{lem:belief_monotone}
$e(\sigma_0,\sigma_1)$ is strictly decreasing in $\sigma_0$ and strictly increasing in $\sigma_1$.
\end{lemma}
\begin{proof}[Proof sketch]
Increasing $\sigma_0$ raises the denominator (more low-quality hiders enter the pool) while leaving the numerator unchanged, so $e$ falls. Increasing $\sigma_1$ raises the numerator and the denominator by the same term $\pi\bar{\theta}$, and since $e < 1$ the numerator effect dominates, so $e$ rises. The strict monotonicity in $\sigma_0$ is the self-correcting force that ensures equilibrium uniqueness in Section~\ref{sec:equilibrium}: as more informed low-quality sellers choose to hide, buyer beliefs worsen, attenuating the benefit of hiding and restoring incentive compatibility.
\end{proof}

\subsection{Equilibrium}
\label{sec:equilibrium}

\subsubsection{Disclosure Incentives}

The seller's disclosure decision turns on a single comparison: whether the price attainable by withholding the audit result, $p^*_{d=0}$, priced off the pool belief $e$, exceeds the price attainable by disclosing true quality $\theta$. A low-quality seller ($\theta=0$) gains from non-disclosure by pooling with uninformed sellers and extracting the associated quality premium; a high-quality seller ($\theta=1$) loses from non-disclosure because pooling with low-quality hiders depresses the belief $e$ below 1. This asymmetry drives the central result.

\begin{proposition}[Disclosure Incentives and Seller Sorting]
\label{prop:premium}
Under the benchmark $\alpha_k=\beta_k$, a seller of type $(K,\theta)$ strictly prefers disclosure if and only if $\theta > e(\sigma_0^*)$. In particular:
\begin{enumerate}[label=(\roman*)]
    \item An informed seller with $\theta=1$ strictly prefers disclosure for any $e(\sigma_0^*)<1$.
    \item An informed seller with $\theta=0$ strictly prefers non-disclosure whenever $e(\sigma_0^*)>0$.
\end{enumerate}
\end{proposition}

\noindent\textit{Proof.} See Appendix~\ref{app:proofs}. $\Box$

Proposition~\ref{prop:premium} reflects a standard adverse-selection asymmetry. For a high-quality seller, disclosure reveals quality exactly, whereas non-disclosure pools her with low-quality hiders; the condition $e(\sigma_0^*)<1$ is satisfied for any non-degenerate pool and so is always binding. For a low-quality seller, concealment allows her to price off the pool belief $e(\sigma_0^*)>0$, which is strictly greater than her true quality $\theta=0$; the gain from pooling is therefore strictly positive whenever uninformed sellers are present. The additional hiding cost $c_1>0$ borne by high-quality concealers reinforces but does not drive this separation under the benchmark $\alpha_k=\beta_k$: the pool-inference argument alone is sufficient to induce full separation of high-quality sellers, with $c_1$ becoming binding only when $\alpha_k < \beta_k$ (the distrust discount case).

\begin{corollary}[High-Quality Sellers Always Disclose]
\label{cor:hq_disclose}
In any equilibrium, $\sigma_1^*=0$.
\end{corollary}

\noindent\textit{Proof.} Follows directly from Proposition~\ref{prop:premium}(i). $\Box$

Corollary~\ref{cor:hq_disclose} is a partial unraveling result: the top of the quality distribution fully separates. An informed high-quality seller gains nothing from ambiguity and is strictly harmed by being pooled with low-quality hiders; she therefore always discloses. The non-disclosing pool is consequently free of high-quality informed sellers in any equilibrium, and the strategic question reduces entirely to the behavior of informed low-quality sellers.

Since $\sigma_1^*=0$, the equilibrium belief simplifies. Defining $\sigma_0\equiv\Pr(d=0\mid K=1,\theta=0)$, the non-disclosing pool consists exclusively of uninformed sellers of both quality types and informed low-quality hiders:
\begin{equation}
    e(\sigma_0) = \frac{(1-\pi)\bar{\theta}}{\pi(1-\bar{\theta})\sigma_0 + (1-\pi)}
    \label{eq:belief_simplified}
\end{equation}
The belief $e(\sigma_0)$ is strictly decreasing in $\sigma_0$: each additional low-quality hider enters the denominator with weight $\pi(1-\bar{\theta})$ while leaving the numerator unchanged, diluting pool quality and depressing buyer beliefs.

\subsubsection{Equilibrium Characterization}

Let $\Delta(e)\equiv\rho\cdot\beta_1\cdot e$ denote the revenue gain from non-disclosure for a low-quality seller facing pool belief $e$. Define the net hiding gain:
\begin{equation}
    \Phi(\sigma_0) \;\equiv\; \Delta\!\left(e(\sigma_0)\right) - c_0 \;=\; \frac{\rho\,\beta_1\,(1-\pi)\bar{\theta}}{\pi(1-\bar{\theta})\sigma_0 + (1-\pi)} - c_0.
    \label{eq:Phi}
\end{equation}
The function $\Phi$ is strictly decreasing in $\sigma_0$: as more informed low-quality sellers conceal, pool quality $e$ deteriorates, the gain from hiding falls, and the incentive to conceal dissipates. This monotone feedback delivers a unique equilibrium without reliance on equilibrium refinements. The three equilibrium regimes follow directly from comparing $\Phi$ to zero at the two boundary values: $\Phi(0)=\Delta(\bar{\theta})-c_0$ measures the net gain when the pool is at its cleanest, and $\Phi(1)=\Delta(e(1))-c_0$ when it is maximally contaminated.

\begin{proposition}[Equilibrium Existence and Uniqueness]
\label{prop:equilibrium}
There exists a unique equilibrium, characterized as follows:
\begin{enumerate}[label=(\roman*)]
    \item If $c_0>\Delta(\bar{\theta})$, the unique equilibrium is \emph{Cost-Deterred Hiding}: $\sigma_0^*=0$ and $e^*=\bar{\theta}$.
    \item If $c_0<\Delta(e(1))$, the unique equilibrium is \emph{Complete Adverse Selection}: $\sigma_0^*=1$ and $e^*=(1-\pi)\bar{\theta}/(1-\pi\bar{\theta})$.
    \item If $\Delta(e(1))\leq c_0\leq\Delta(\bar{\theta})$, the unique equilibrium is \emph{Partial Pooling}: $\sigma_0^*\in(0,1)$ is the unique solution to $\Delta(e(\sigma_0))=c_0$.
\end{enumerate}
\end{proposition}

\noindent\textit{Proof.} See Appendix~\ref{app:proofs}. $\Box$

\begin{table}[h]
\centering
\caption{Equilibrium Regimes}
\label{tab:equilibria}
\small\renewcommand{\arraystretch}{1.5}
\begin{tabular}{lcccc}
\toprule\toprule
\textbf{Equilibrium} & \textbf{Condition} & \textbf{Strategy $\sigma_0^*$} & \textbf{Belief $e^*$} & \textbf{Premium $p_0^*-p_1^*$} \\
\midrule
Cost-Deterred Hiding       & $c_0>\Delta(\bar{\theta})$                     & $0$     & $\bar{\theta}$                                          & $\rho\beta_1\bar{\theta}$ \\
Complete Adverse Selection & $c_0<\Delta(e(1))$                             & $1$     & $\dfrac{(1-\pi)\bar{\theta}}{1-\pi\bar{\theta}}$        & $\rho\beta_1 e^*$ \\[6pt]
Partial Pooling            & $\Delta(e(1))\leq c_0\leq\Delta(\bar{\theta})$ & $(0,1)$ & Solves $\Phi(\sigma_0^*)=0$                             & $\rho\beta_1 e(\sigma_0^*)$ \\
\bottomrule\bottomrule
\end{tabular}
\end{table}

Neither corner solution fits the data: universal concealment would require the non-disclosure premium to exceed $c_0$ even against the most contaminated pool, while universal disclosure would require concealment to be unprofitable even against the cleanest. What we observe instead is an interior equilibrium in which some low-quality sellers disclose and the remainder pool with uninformed non-disclosers. Substituting~\eqref{eq:Phi} into $\Phi(\sigma_0^*)=0$ yields the closed-form hiding rate:
\begin{equation}
    \sigma_0^* = \frac{1}{\pi(1-\bar{\theta})}\left[\frac{\rho\,\beta_1\,(1-\pi)\bar{\theta}}{c_0} - (1-\pi)\right].
    \label{eq:pp_sigma}
\end{equation}

We derive the non-disclosure premium as follows:
\begin{equation}
    p_0^*-p_1^* \;=\;
    \begin{cases}
        \rho\beta_1\bar{\theta}
            & \text{Cost-Deterred Hiding }(\sigma_0^*=0), \\[4pt]
        c_0
            & \text{Partial Pooling }(\sigma_0^*\in(0,1)), \\[4pt]
        \rho\beta_1\dfrac{(1-\pi)\bar{\theta}}{1-\pi\bar{\theta}}
            & \text{Complete Adverse Selection }(\sigma_0^*=1),
    \end{cases}
    \quad \text{with}\quad \rho\beta_1\frac{(1-\pi)\bar{\theta}}{1-\pi\bar{\theta}} < c_0 < \rho\beta_1\bar{\theta}.
    \label{eq:premium}
\end{equation}

The three equilibria characterized in Proposition~\ref{prop:equilibrium} arise from a single comparison: the hiding gain $\Delta(e)=\rho\beta_1 e$ evaluated at the two extremes of pool contamination against the structural hiding cost $c_0$. When the gain is dominated by $c_0$ even at the cleanest possible pool ($\sigma_0=0$, $e=\bar{\theta}$), no informed seller hides; when it survives the most contaminated pool ($\sigma_0=1$, $e=(1-\pi)\bar{\theta}/(1-\pi\bar{\theta})$), every informed seller hides; in between, an interior fraction hides and the system self-equilibrates. The non-disclosure premium inherits a strict ordering across regimes,
\begin{equation*}
    \underbrace{\rho\beta_1 \bar{\theta}}_{\text{Cost-Deterred Hiding}}
    \;>\;
    \underbrace{c_0}_{\text{Partial Pooling}}
    \;>\;
    \underbrace{\rho\beta_1\,\dfrac{(1-\pi)\bar{\theta}}{1-\pi\bar{\theta}}}_{\text{Complete Adverse Selection}},
\end{equation*}
with the interior value depending neither on $\pi$ nor on $\rho$.

In the \emph{Cost-Deterred Hiding} regime, the hiding cost is prohibitive even against the cleanest conceivable non-disclosing pool, the one populated entirely by uninformed sellers drawn from the unconditional quality distribution. No informed low-quality seller finds it profitable to suppress her type, so $\sigma_0^*=0$ and the non-disclosing pool is uncontaminated by strategic behavior: $e^*=\bar{\theta}$. The premium $\rho\beta_1\bar{\theta}$ reflects the full unconditional quality belief, scaled by the share of EPC-sensitive buyers and their marginal valuation. Relative to \citet{milgrom_good_1981} and \citet{grossman_informational_1981}, this regime delivers a form of deterrence (informed bad types refrain from concealment) but through a fundamentally different mechanism: unraveling in those models is sustained by a belief-penalty channel, in which silence triggers the worst possible inference, whereas here deterrence runs through the cost channel $c_0$ because uninformed sellers prevent silence from carrying the worst-type inference. This is the only regime in which the non-disclosing pool is informationally clean.

In the \emph{Complete Adverse Selection} regime, the hiding gain survives the most contaminated pool the model can sustain. Every informed low-quality seller hides, $\sigma_0^*=1$, and buyer beliefs about the non-disclosing pool collapse to the floor $e^*=(1-\pi)\bar{\theta}/(1-\pi\bar{\theta})$, strictly below $\bar{\theta}$. The non-disclosure premium accordingly hits its lower bound. This lower bound is nonetheless strictly positive: the uninformed seller mass dilutes the adverse inference and prevents the complete belief collapse that Milgrom--Grossman unraveling would imply. The mechanism is exactly the one identified by \citet{dye_disclosure_1985}, in which a residual silent population that is not strategically motivated keeps silence partially informative, and it operates here even when strategic hiders saturate.

In the \emph{Partial Pooling} regime, $\Phi(\sigma_0)$ is positive at $\sigma_0=0$ and negative at $\sigma_0=1$. Strict monotonicity yields a unique interior crossing $\sigma_0^*\in(0,1)$ at which the marginal hider is exactly indifferent: $\Delta(e(\sigma_0^*))=c_0$. The underlying mechanism is a feedback loop. As more informed low-quality sellers join the non-disclosing pool, $e(\sigma_0)$ falls, perceived pool quality deteriorates, and the gain $\Delta(e)$ from hiding shrinks. Entry into the pool stops precisely when this gain has been driven down to the hiding cost. The fraction $\sigma_0^*$ of low-quality types hides and pools with the uninformed; the remaining $1-\sigma_0^*$ disclose, accepting the low-quality price rather than continuing to contaminate the pool.

The structural payoff of this regime is that the equilibrium premium equals exactly $c_0$, invariant to $\pi$ and $\rho$ throughout the interior. The reason is the indifference condition itself: in equilibrium, the price difference between the non-disclosing pool and the disclosed low-quality outcome must equal the hiding cost, otherwise the marginal seller would strictly prefer one action over the other. To see the invariance, consider a rise in $\pi$. A larger informed share would, holding $\sigma_0^*$ fixed, mechanically contaminate the pool and depress $e$. Equilibrium restores $\Phi=0$ by reducing $\sigma_0^*$: fewer hiders produce a cleaner pool, $e$ rebounds, and the premium returns to $c_0$. An analogous adjustment absorbs changes in $\rho$. The non-disclosure premium therefore functions as a revealed hiding cost rather than a free function of market composition, and the observable price gap between non-disclosing and disclosing low-quality sellers directly identifies $c_0$. This is the only regime in which $\pi$, $\rho$, and $c_0$ are jointly identified.

\subsubsection{Comparative Statics}

\begin{proposition}[Comparative Statics]
\label{prop:comparative_statics}
The following comparative statics hold:
\begin{enumerate}[label=(\roman*)]
    \item In the \emph{Complete Adverse Selection} equilibrium: $\partial(p_0^*-p_1^*)/\partial\pi < 0$ and $\partial(p_0^*-p_1^*)/\partial\rho > 0$.
    \item In the \emph{Partial Pooling} equilibrium: $\partial\sigma_0^*/\partial c_0 < 0$, $\partial\sigma_0^*/\partial\rho > 0$, and $\partial\sigma_0^*/\partial\pi < 0$.
\end{enumerate}
\end{proposition}

\noindent\textit{Proof.} See Appendix~\ref{app:proofs}. $\Box$

\begin{itemize}
    \item Part~(i) reflects two distinct transmission channels. The effect of $\pi$ operates through the \textit{belief channel}: as the share of informed sellers rises, non-disclosure becomes more informative of low quality, $e^*$ falls, and the gain from concealment dissipates. The effect of $\rho$ operates through the \textit{revenue channel}: a larger share of EPC-sensitive buyers widens the price spread between quality levels, raising the value of concealment.
    \item Part~(ii) follows from the implicit function theorem applied to $\Phi(\sigma_0^*)=0$. Since $\Phi'(\sigma_0^*)<0$, one obtains $\partial\sigma_0^*/\partial c_0 < 0$ directly. The effect $\partial\sigma_0^*/\partial\rho > 0$ follows because a higher $\rho$ increases $\Delta$ at any fixed $e$, shifting $\Phi$ upward and raising the interior equilibrium. The effect of $\pi$ operates through the belief function: a higher $\pi$ increases the contamination weight $\pi(1-\bar{\theta})$ in the denominator of $e(\sigma_0)$, reducing $\Delta(e(\sigma_0))$ and thereby shifting $\Phi$ downward, so $\partial\sigma_0^*/\partial\pi < 0$. The parameters $\pi$ and $c_0$ are therefore \textit{policy substitutes}: each reduces strategic hiding through a distinct channel ($\pi$ through beliefs, $c_0$ through costs), and their marginal effects on $\sigma_0^*$ are both negative and quantifiable from the estimated model.
\end{itemize}

\begin{corollary}[Effect of $\rho$ on the non-disclosure premium]
\label{cor:rho}
\begin{enumerate}[label=(\roman*)]
    \item The non-disclosure premium $\Delta = p_0^* - p_1^*$ is weakly increasing in $\rho$: strictly so in the Cost-Deterred Hiding and Complete Adverse Selection regimes ($\partial\Delta/\partial\rho > 0$), and zero in the Partial Pooling regime ($\partial\Delta/\partial\rho = 0$, with the premium pinned at $c_0$).
    \item The parameter interval sustaining the Partial Pooling regime, $[\rho\,\beta_1\,e(1),\,\rho\,\beta_1\,\bar{\theta}]$, is weakly expanding in $\rho$: its width $\rho\,\beta_1[\bar{\theta}-e(1)]$ is strictly increasing in $\rho$ whenever $\pi>0$.
\end{enumerate}
\end{corollary}

\noindent\textit{Proof.} See Appendix~\ref{app:proofs}. $\Box$

\noindent Economically, a larger share of buyers who price the certificate raises what silence
is worth wherever the hiding rate is at a corner, but inside the partial-pooling regime it does
not: there the premium is pinned to the concealment cost, and an increase in $\rho$ is absorbed
by sellers hiding more rather than by the price gap widening. Part~(ii) adds that higher $\rho$
widens the parameter window over which partial pooling survives. Empirically this predicts that
across partial-pooling markets the non-disclosure premium should be uncorrelated with
$\hat\rho_m$ while the hiding rate rises with it, and that markets with higher $\hat\rho_m$
should be less likely to appear in a corner regime --- a pattern we can check directly against
the regime classification reported in Section~\ref{sec:results}.


\begin{proposition}[Equilibrium effect of $\pi$]
\label{prop:pi}
As $\pi$ increases, the lower boundary $\Delta(e(1))$ of the Partial Pooling interval falls ($\partial e(1)/\partial\pi < 0$), so the condition for Complete Adverse Selection ($c_0 < \Delta(e(1))$) becomes harder to satisfy and the equilibrium moves toward the Cost-Deterred Hiding regime. For any fixed hiding cost $c_0$, a sufficiently large increase in $\pi$ drives the equilibrium out of the Partial Pooling region toward full unraveling.
\end{proposition}

\noindent\textit{Proof.} See Appendix~\ref{app:proofs}. $\Box$

\noindent In words: better registry coverage makes silence more informative, because a larger
fraction of non-disclosers are hiding by choice rather than by constraint. Buyers therefore
discount silence more heavily, concealment becomes less attractive, and the market moves toward
disclosure. The testable content is a cross-market comparative static: markets with higher
$\hat\pi_m$ should display lower hiding rates $\hat\sigma_{0,m}$ and should be
disproportionately represented among the markets classified as fully unravelling in
Section~\ref{sec:results}.


\begin{proposition}[Equilibrium effect of $c_0$]
\label{prop:c0}
An increase in $c_0$ shifts $\Phi(\sigma_0) = \Delta(e(\sigma_0)) - c_0$ uniformly downward. The upper boundary condition for Partial Pooling ($c_0 \leq \Delta(\bar{\theta})$) becomes harder to satisfy and the condition for Cost-Deterred Hiding ($c_0 > \Delta(\bar{\theta})$) becomes easier, so higher hiding costs push the equilibrium toward the Cost-Deterred Hiding regime. In the interior of the Partial Pooling regime, $\partial\sigma_0^*/\partial c_0 < 0$: higher costs reduce the equilibrium hiding rate.
\end{proposition}

\noindent\textit{Proof.} See Appendix~\ref{app:proofs}. $\Box$

The premium in~\eqref{eq:premium} generates a transfer from EPC-sensitive buyers to the non-disclosing pool, but it does not accrue symmetrically within that pool. The informed low-quality hider receives the premium but nets zero: the indifference condition $\Delta(e(\sigma_0^*))=c_0$ ensures that the gross revenue gain from concealment is exactly consumed by the hiding cost. The uninformed seller, who bears no hiding cost, is the residual claimant and pockets the full premium $c_0$ as a pure rent. The welfare cost of the equilibrium therefore has two distinct components. First, an allocative distortion: EPC-sensitive buyers pay a positive premium for assets whose true energy performance falls below the pooled belief $e(\sigma_0^*)$, a misallocation that would be eliminated if the concealed information were transmitted. Second, a deadweight loss: the hiding cost $c_0$ is incurred by informed sellers at listing time without any counterpart in buyer valuations. This is a social loss only to the extent that $c_0$ represents real resources destroyed (legal fees, prolonged negotiation, administrative burden); any portion that is a transfer, such as the expected value of a fine accruing to the state, is a redistribution rather than a loss. We do not identify this split, so the deadweight-loss component is an upper bound. A formal efficiency characterization of these two components and the optimal subsidy that addresses them is developed in Section~\ref{sec:policy}.

\subsection{Identification}
\label{sec:identification}

The two structural objects to identify are $\rho_m$, the market-specific share of EPC-sensitive buyers, and $c_{0,m}$, the seller hiding cost. The structural equation $\Delta_m = \rho_m \cdot \beta_{1,m} \cdot e^*_m$ relates these to observables: $\Delta_m$ and $\beta_{1,m}$ are equilibrium price gaps, and $e^*_m$ is directly observed from the ADEME registry. This might appear to make $\rho_m$ point-identified from a single cross-section via $\rho_m = \Delta_m / (\beta_{1,m} \cdot e^*_m)$. The problem is that cross-sectional price gaps reflect \emph{level} differences across markets, not the demand-induced component we need. Any market characteristic that simultaneously shapes the composition of the undisclosed pool and is correlated with buyer EPC-sensitivity would bias a cross-sectional estimate of $\rho_m$. Purging this confound requires variation that isolates the \emph{change} in price gaps attributable to buyers becoming EPC-sensitive, not their levels.

The July 2021 DPE reform provides this variation, and the two-step approach uses it as follows. In the first step, a difference-in-differences estimator compares price gaps in each market before and after the reform, stripping out market-specific level differences and isolating the reform-induced shift in buyer demand locally. This yields $\hat\rho_m = \hat\Delta_m / (\hat\beta_{1,m} \cdot \hat{e}^*_m)$, where the hat quantities are DiD estimates (post minus pre period means). In the second step, the hiding cost $c_{0,m}$ is recovered from the equilibrium hiding rate $\hat\sigma_{0,m}$ via a closed-form inversion of the seller indifference condition (given the equilibrium characterized in Section~\ref{sec:model}). No additional external variation is required: $\hat\sigma_{0,m}$ is directly observed as the fraction of informed low-quality sellers who withhold their EPC, and the inversion is exact. Together, three observables $(\hat\beta_{1,m},\,\hat\Delta_m,\,\hat\sigma_{0,m})$ and one registry-based input $(\hat{e}^*_m)$ map to two structural parameters $(\hat\rho_m,\,\hat{c}_{0,m})$, market by market.

\subsubsection{The July 2021 Reform as a Natural Experiment}

The change in the legal status of the DPE described in Section~\ref{ss:institutions} is the source of identification: it created a precise, legally documented change in the value of a disclosed EPC score to buyers, with a known implementation date of July 1, 2021. Crucially, the reform operated on the \emph{demand} side --- the weight a disclosed rating carries in buyers' valuations and legal recourse --- and not on the seller's cost of withholding. The obligation to provide a DPE at sale, and the compliance and liability exposure attached to coming to market without the required diagnostic, were already in force under loi ALUR (Article L271-4 CCH) well before 2021. These are precisely the components of the concealment cost $c_{0,m}$. What July 2021 changed was not whether a seller must disclose, nor the cost she bears for withholding, but how much a \emph{disclosed} rating now matters to the buyer. The reform therefore shifts $\rho_m$ while leaving $c_{0,m}$, governed by the pre-existing disclosure mandate, unchanged across the reform date. This is what licenses reading the DiD change in the premium as the demand response that identifies $\rho_m$, holding the hiding cost fixed. Whether the supply side of the market responded concurrently is a maintained assumption we cannot directly test with the data at hand; we return to it below.\footnote{An alternative strategy based on a Bartik instrument (interacting wholesale gas price variation with the pre-reform share of low-efficiency properties per market) fails at the first stage ($F$-statistic: 1.0 to 1.6). The French \emph{bouclier tarifaire} (energy price shield, October 2021 to 2023) decoupled household energy costs from the TTF gas price spike driven by the Russian invasion of Ukraine, removing the demand-side variation the instrument was designed to exploit.}

The identifying assumption is the standard parallel-trends condition: in the absence of the reform, the cross-group price gaps $\beta_{1,m,t}$ (energy-efficient vs.\ energy-intensive disclosed properties) and $\Delta_{m,t}$ (undisclosed vs.\ energy-intensive) would have continued on their pre-reform trajectory. The DiD does not require that buyers were entirely inattentive to EPC scores before July 2021; it requires only that their relative pricing of disclosure groups was not already trending. Any post-reform shift in the gaps, relative to that pre-reform baseline, is then attributable to the change in buyer incentives induced by legal opposability.

The pre-trend evidence supports this assumption directly, and more strongly than a single test would. Both $\beta_{1,m,t}$ and $\Delta_{m,t}$ are statistically flat across all markets in the 33 months before the reform (Figure~\ref{fig:ATT}). These two gaps measure distinct economic objects: $\beta_{1,m,t}$ captures buyers' willingness to pay for verified high quality, while $\Delta_{m,t}$ reflects the informational wedge created by strategic non-disclosure. They respond to different primitives of the model. A supply-side confound (a shift in $\pi_m$, $\bar\theta_m$, or seller composition at the reform date) would generically move these two gaps by different amounts and in potentially different directions, since the mechanisms through which supply composition enters each gap differ. The fact that both gaps are jointly flat before July 2021 and both shift upward simultaneously after the reform, at the same markets and on the same calendar date, is difficult to reconcile with a supply-side story. It is exactly what the demand-side interpretation predicts: legal opposability increases the return to EPC quality for sensitive buyers, raising both the price of verified high quality ($\beta_{1,m,t}$) and the information premium extracted by non-disclosers ($\Delta_{m,t}$) at the same moment.

We acknowledge that supply-side stability (the assumption that $\pi_m$, $\bar\theta_m$, and $c_{0,m}$ did not shift discontinuously at the reform) is maintained rather than directly tested. Our listing data begin in 2021Q3, so pre-reform ADEME match rates are not observable. The timing of audit decisions provides some reassurance: certified DPE audits are commissioned and completed before listing, with a gap of weeks to months, so the audit stock at any listing date reflects decisions made prior to that date. But this is a maintained assumption, and we treat it as such. The consequence of a supply-side shift at the reform is bounded: any resulting bias in $\hat\rho_m$ is upward, reinforcing the interpretation of $\hat\rho_m$ as an upper bound rather than a point estimate.\footnote{Because the disclosure obligation and its associated liability predate the reform (loi ALUR, Article L271-4 CCH), we treat $c_{0,m}$ as governed by that long-standing mandate and unchanged across July 2021. In the conservative case that opposability nonetheless nudged the perceived cost of withholding upward, $\hat\Delta_m$ would absorb both the demand response and the cost change, causing $\hat\rho_m = \hat\Delta_m / (\hat\beta_{1,m}\cdot\hat{e}^*_m)$ to overstate true buyer sensitivity. This is the same direction of bias as from supply-side instability: both reinforce the upper-bound interpretation.}

\subsubsection{DiD Identification of $\rho_m$}

The identification uses a panel of monthly hedonic price indices constructed by Chapter~1 from listing-level data covering all three disclosure groups. For each market $m$, group $g\in\{H,\,L,\,\emptyset\}$, and month $t$, the index $\bar{p}_{m,t}^g$ is the predicted log price per m$^2$ for a reference property conditional on hedonic controls $X_i$. The $H$ index covers energy-efficient disclosed properties (EPC grades A to D, $d=1$, $K=1$); the $L$ index covers energy-intensive disclosed properties (grades E to G, $d=1$, $K=1$); and the $\emptyset$ index covers all non-disclosing listings ($d=0$), a pooled group that blends strategic hiders ($K=1$) and genuinely uninformed sellers ($K=0$). By conditioning on observables, the cross-group gap $\bar{p}_{m,t}^g - \bar{p}_{m,t}^{g'}$ isolates the EPC-attributable price wedge net of compositional differences between groups.

From these indices, two time-varying log price gaps are defined:
\begin{align}
\beta_{1,m,t} &= \bar p_{m,t}^{H} - \bar p_{m,t}^{L} \quad \text{(energy-efficiency premium)} \\
\Delta_{m,t}  &= \bar p_{m,t}^{\emptyset} - \bar p_{m,t}^{L} \quad \text{(non-disclosure premium)}
\end{align}
The reform-induced shifts in each gap are estimated by differencing period means:
\begin{align}
\hat\beta_{1,m} &= \overline{\beta_{1,m,t}}^{\,\text{post}} - \overline{\beta_{1,m,t}}^{\,\text{pre}}
  \label{eq:did_beta1} \\
\hat\Delta_m    &= \overline{\Delta_{m,t}}^{\,\text{post}} - \overline{\Delta_{m,t}}^{\,\text{pre}}
  \label{eq:did_delta}
\end{align}
where the pre-period spans October~2018 through June~2021 (33~months) and the post-period spans July~2021 through July~2024 (37~months). A volume filter requiring at least 20 average monthly transactions per disclosure group per market retains 40 apartment markets and 100 house markets; markets with $\hat\beta_{1,m}\leq 0$ are excluded from structural inversion.

To validate the parallel pre-trends assumption and characterise the dynamic profile of the reform, the period-mean estimates are complemented with a month-by-month event study. For each market $m$ and relative month $k$ (where $k=0$ is July 2021), the normalised gaps
\begin{align}
\tilde\beta_{1,m,k} &= \beta_{1,m,t(k)} - \beta_{1,m,t(-1)}, \qquad
\tilde\Delta_{m,k}  = \Delta_{m,t(k)}  - \Delta_{m,t(-1)}
\end{align}
are set to zero at $k=-1$ (the last pre-reform month) by construction. Figure~\ref{fig:ATT} plots the cross-market means with 95\% confidence intervals over $k\in[-33,\,+36]$. The pre-reform coefficients ($k<0$) are statistically indistinguishable from zero for both gaps, confirming parallel pre-trends. The post-reform coefficients increase gradually from $k=0$, consistent with progressive capitalisation of the legally opposable DPE. The period-mean estimates in Equations~\eqref{eq:did_beta1} and \eqref{eq:did_delta} average these relative-month profiles over the respective windows.

Given these estimates, the structural equation $\Delta_m = \rho_m\cdot\beta_{1,m}\cdot e^*_m$ (Equation~\eqref{eq:premium}) implies a closed-form estimator for buyer sensitivity. Since the reform shifts $\rho_m$ while $e^*_m$ is a supply-side object measured directly from the ADEME registry (Section~\ref{sec:beliefs} below), the DiD isolates the post-reform demand response:
\begin{equation}
    \hat\rho_m = \frac{\hat\Delta_m}{\hat\beta_{1,m} \cdot \hat{e}^*_m}
    \label{eq:rho_estimator}
\end{equation}
The constraint $\hat\rho_m \leq 1$ binds in the majority of markets. For apartments, the unconstrained ratio exceeds unity in all 40 markets (cross-market mean of approximately 7.7); for houses, 92 of 107 markets with $\hat\beta_{1,m}>0$ are clipped (cross-market mean ratio of approximately 0.58). Two compounding sources drive the excess: composition bias in $\hat\Delta_m$, because the non-disclosure price index pools strategic hiders with genuinely uninformed sellers and inflates the numerator; and thin-cell attenuation of $\hat\beta_{1,m}$, because sparse disclosure-group cells compress the denominator. These effects are quantitatively consistent with the magnitudes observed, but they preclude point identification of $\rho_m$ in most markets. In clipped markets, $\rho_m = 1$ is maintained as a working assumption, and $\hat{c}_{0,m}$ is interpreted as an upper bound: since the indifference condition $c_{0,m} = \rho_m\cdot\beta_{1,m}\cdot\hat{e}^*_m(\hat\sigma_{0,m})$ is strictly increasing in $\rho_m$, setting $\rho_m = 1$ overstates $c_{0,m}$ relative to any interior value. Point identification is achieved in 8 house markets where the inversion yields an interior solution; these are characterised separately in the Results.

\subsubsection{Measuring Equilibrium Beliefs}
\label{sec:beliefs}

Computing $\hat\rho_m$ and $\hat{c}_{0,m}$ requires $\hat{e}^*_m$, the expected quality buyers assign to a non-disclosing listing. In standard signalling models this belief must be jointly estimated with the structural parameters. Here it is directly observable: the ADEME registry provides the researcher with the same information structure that $e^*_m$ encodes for buyers.

For each market $m$, the four-group classification of Section~\ref{sec:data} yields three sample frequencies: $\hat\pi_m$ (the share of listings with $K=1$, the ADEME match rate); $\hat{\bar\theta}_m$ (the share of informed sellers with high quality, $\theta=1$, EPC $\leq$ D); and $\hat\sigma_{0,m}$ (the share of informed low-quality sellers, $K=1$, $\theta=0$, who hide by setting $d=0$). Substituting into the belief function (Equation~\eqref{eq:belief_simplified}, at $\sigma_1^*=0$ per Corollary~\ref{cor:hq_disclose}):
\begin{equation}
    \hat{e}^*_m = \frac{(1-\hat\pi_m)\,\hat{\bar\theta}_m}{\hat\pi_m\,(1-\hat{\bar\theta}_m)\,\hat\sigma_{0,m} + (1-\hat\pi_m)}
    \label{eq:estar}
\end{equation}
The numerator is the uninformed high-quality mass in the non-disclosing pool; the denominator is the total non-disclosing mass. No model estimation is involved: $\hat{e}^*_m$ is a function of directly observed sample proportions.

\subsubsection{Closed-Form Identification of $c_{0,m}$}

One might expect $\hat{c}_{0,m} = \hat\Delta_m$, since the equilibrium indifference condition gives $c_{0,m} = \Delta_m$ for a marginal hider. This logic fails because the ``undisclosed'' category in the price index pools \emph{two heterogeneous groups}: strategic hiders ($K=1$, $\theta=0$, $d=0$), who are the economically relevant group, and genuinely uninformed sellers ($K=0$, $d=0$), who constitute approximately 85\% of all non-disclosers. $\hat\Delta_m$ is therefore a weighted average of the strategic hider premium and the uninformed-seller premium, an upper bound on $c_{0,m}$, not a point estimate.\footnote{A listing-level DiD targeting $K=1$ non-disclosers directly is infeasible: the listing data begin in 2021, providing no pre-reform baseline at the listing level.}

The route to $c_{0,m}$ uses the equilibrium hiding rate $\hat\sigma_{0,m}$ among informed low-quality sellers ($K=1$, $\theta=0$). This group is identified precisely by the ADEME match and is immune to contamination from uninformed sellers. In the Partial Pooling regime, the equilibrium indifference condition
\begin{equation}
    \Delta\!\left(e^*(\hat\sigma_{0,m})\,;\,\hat\rho_m\right) = c_{0,m}
    \label{eq:indifference}
\end{equation}
has a unique solution in $c_{0,m}$ for any observed $\hat\sigma_{0,m}\in(0,1)$. Given the first-stage estimate $\hat\rho_m$, this yields the closed-form estimator:
\begin{equation}
    \hat{c}_{0,m} = \hat\rho_m \cdot \hat\beta_{1,m} \cdot \hat{e}^*_m(\hat\sigma_{0,m})
    \label{eq:c0_estimator}
\end{equation}
where $\hat{e}^*_m(\hat\sigma_{0,m})$ is the belief evaluated at the observed mixing rate. The estimator is exact: one free parameter ($c_{0,m}$) is matched to one moment ($\hat\sigma_{0,m}$) via the equilibrium condition. No numerical optimisation is required.

For house markets, the data provide two distinct upper bounds on $c_{0,m}$. The first is the equilibrium inversion under the maintained $\hat\rho_m = 1$ assumption:
\begin{equation}
    \hat{c}_{0,m}^{\,\text{inv}} = \hat\beta_{1,m} \cdot \hat{e}^*_m(\hat\sigma_{0,m})
    \label{eq:c0_interval}
\end{equation}
The second is $\hat\Delta_m$, the DiD non-disclosure premium, which is an upper bound on $c_{0,m}$ because uninformed sellers (approximately 75\% of non-disclosers for houses) inflate the undisclosed pool's average price above what strategic hiders alone would generate. The cross-market means are $\hat{c}_{0,m}^{\,\text{inv}} = 0.018$ log points and $\hat\Delta_m = 0.034$ log points. Internal consistency requires $\hat{c}_{0,m}^{\,\text{inv}} \leq \hat\Delta_m$, which holds in all but a handful of markets. The tighter bound $\hat{c}_{0,m}^{\,\text{inv}} = 0.018$ is the preferred estimate and is used in the policy counterfactuals of Section~\ref{sec:policy}.

The structural estimates use listing-level data from 2022 to 2024.\footnote{The year 2021 is excluded from the listing-level analysis because the market had not yet fully repriced the reform: the pooled hedonic estimate of $\hat\beta_{1}$ for apartments in H2~2021 is approximately $-0.012$, indicating the EPC quality premium had not yet turned positive, which renders the structural equation $\Delta_m = \rho_m\cdot\beta_{1,m}\cdot e^*_m$ unreliable for that period.} The estimation sample comprises 40 apartment markets and 100 house markets.

Bootstrap standard errors are computed with $B=500$ within-market resamples. In each draw, listings in market $m$ are resampled with replacement and all market-level moments $(\hat\pi_m, \hat{\bar\theta}_m, \hat\sigma_{0,m})$ are recomputed. The first-stage estimate $\hat\rho_m$ is held fixed at its DiD value across bootstrap iterations; only $\hat\sigma_{0,m}$ is resampled to propagate sampling uncertainty into $\hat{c}_{0,m}$ via Equation~\eqref{eq:c0_estimator}. The 500 replications yield percentile-based 95\% confidence intervals.

\subsection{Estimates}
\label{sec:results}

\subsubsection{Main Estimates}

Table~\ref{tab:did_results} reports market-level estimates for both property types. Panel A covers the DiD estimates from the price index panel; Panel B covers the two-step structural estimates; Panel C covers supply-side observables from the ADEME match.

\begin{table}[H]
\centering
\caption{Market-level estimates by property type}
\label{tab:did_results}
\small\renewcommand{\arraystretch}{1.3}
\begin{tabular*}{\textwidth}{@{\extracolsep{\fill}}l*{6}{c}@{}}
\toprule\toprule
 & \multicolumn{3}{c}{\textbf{Apartments}} & \multicolumn{3}{c}{\textbf{Houses}} \\
\cmidrule(lr){2-4}\cmidrule(lr){5-7}
\textbf{Parameter} & \textbf{Mean} & \textbf{Std} & \textbf{N} & \textbf{Mean} & \textbf{Std} & \textbf{N} \\
\midrule
\multicolumn{7}{l}{\textit{Panel A: DiD estimates from price index panel}} \\[2pt]
$\hat\beta_{1,m}$ (energy-efficiency premium) & 1.7\% & 3.5 & 49  & 5.1\% & 4.0 & 104 \\
$\hat\Delta_m$ (non-disclosure premium, UB)   & 6.9\% & 6.3 & 49  & 3.4\% & 4.8 & 104 \\
\midrule
\multicolumn{7}{l}{\textit{Panel B: Two-step structural estimates, pooled 2022 to 2024}} \\[2pt]
$\hat\rho_m$ (buyer EPC-sensitivity)          & 90.9\% & 28.6 & 45 & 68.1\% & 42.0 & 97 \\
$\hat c_{0,m}$ (hiding cost)                  &  3.4\% &  2.9 & 45 &  4.5\% &  3.7 & 97 \\
\midrule
\multicolumn{7}{l}{\textit{Panel C: Supply-side observables from ADEME match}} \\[2pt]
$\hat\pi_m$ (registry-matched share)          & 50.7\% &  7.4 & 45 & 41.9\% &  4.4 & 97 \\
$\hat{\bar\theta}_m$ (mean quality $K{=}1$)   & 62.5\% &  8.8 & 45 & 60.5\% & 10.4 & 97 \\
$\hat\sigma_{0,m}$ (informed hiding rate)     & 14.0\% & 10.5 & 45 &  5.7\% &  6.5 & 97 \\
\bottomrule\bottomrule
\end{tabular*}
\vspace{0.5em}
\begin{minipage}{\linewidth}
\footnotesize\textit{Notes.} Panel A: market-level DiD estimates from the price index panel (pre: Oct 2018 to Jun 2021; post: Jul 2021 to Jul 2024), volume-filtered at $\geq 20$ transactions/month/EPC category. $N=43$ apartment markets after excluding 25 with $\hat\beta_{1,m}\leq 0$; $N=107$ house markets after excluding 6. $\hat\Delta_m$ is an upper bound for $c_{0,m}$ because the undisclosed category pools strategic hiders and uninformed sellers. Panel B: two-step structural estimates using 2022 to 2024 listing data. $\hat\rho_m = \hat\Delta_m / (\hat\beta_{1,m} \cdot \hat{e}^*_m)$ from closed-form DiD inversion, clipped to $[0,1]$; $\hat{c}_{0,m}$ is recovered by equilibrium inversion from $\hat\sigma_{0,m}$ with $\hat\rho_m$ maintained at the clipped value (upper bound when $\hat\rho_m = 1$, since $c_{0,m} = \rho_m\cdot\hat\beta_{1,m}\cdot\hat{e}^*_m$ is increasing in $\rho_m$). Panel C: supply-side moments computed as sample frequencies from the ADEME match. Means in Panel~A and the $\hat c_{0,m}$ row of Panel~B are percentage changes (computed as $e^\delta - 1$, expressed in percent). Share parameters in Panels~B and~C ($\hat\rho_m$, $\hat\pi_m$, $\hat{\bar\theta}_m$, $\hat\sigma_{0,m}$) are in percentage points. Standard-deviation columns are in the same units as the mean (percentage points).
\end{minipage}
\end{table}

\vspace{-1em}

Table~\ref{tab:did_results} yields a coherent and interpretable picture across all three panels.
In Panel~A, the energy-efficiency premium $\hat\beta_{1,m}$ is approximately 1.7\% for
apartments and 5.1\% for houses: disclosing an energy-efficient EPC score commands that price advantage
over disclosing an energy-intensive score, after the July~2021 reform. The non-disclosure premium
$\hat\Delta_m$ (6.9\% for apartments and 3.4\% for houses) measures how much more
non-disclosers earn over low-quality disclosed properties compared with the pre-reform
baseline; because the undisclosed pool commingles strategic hiders with genuinely uninformed
sellers, this figure is an upper bound on the true concealment cost rather than a direct
estimate of it. Panel~B translates these reduced-form inputs into seller-side parameters: $\hat\rho_m = 90.9\%$ for apartments and 68.1\% for
houses indicate that nearly all buyers in the average market price energy performance into
their offers, while the inverted hiding cost $\hat{c}_{0,m}$ settles at 3.4\% of property value for
apartments and 4.5\% for houses --- an economically substantial margin, and one of the same
order in both segments despite how differently they are traded. Panel~C
illuminates the supply side: just over half of apartment sellers (50.7\%) and 41.9\% of house
sellers carry a post-2021 registry record at listing ($\hat\pi_m$), and among those, roughly
six in ten hold a property rated A--D ($\hat{\bar\theta}_m = 62.5\%$ and $60.5\%$).
The most striking contrast across property types lies in the strategic-hiding rate
$\hat\sigma_{0,m}$: 14.0\% of informed apartment sellers who hold an E/F/G rating choose to
withhold it, versus 5.7\% for houses, a difference that mirrors the three-fold gap in
$\hat\beta_{1,m}$, since a larger forgone quality premium makes concealment a costlier choice
for house sellers.

\subsubsection{Structural Estimates and Economic Interpretation}

The first-stage estimator $\hat\rho_m = \hat\Delta_m / (\hat\beta_{1,m} \cdot \hat{e}^*_m)$, clipped to $[0,1]$, yields means of 0.909 for apartments and 0.681 for houses. These averages, however, conceal a pattern of boundary saturation. For apartments, 40 of the 42 partial-pooling markets reach the $\hat\rho_m = 1$ upper boundary; only the 3 full-unraveling markets pull the mean below unity. For houses, 54 of 76 partial-pooling markets reach the boundary, with 21 full-unraveling markets contributing $\hat\rho_m = 0$ and 22 markets achieving interior identification. The boundary reflects a structural feature of the data: the non-disclosure premium $\hat\Delta_m$ typically exceeds $\hat\beta_{1,m} \cdot \hat{e}^*_m$, so the data are consistent with full buyer EPC-sensitivity but cannot pin down exactly how close to one $\rho_m$ truly is. We maintain $\hat\rho_m = 1$ as the working assumption for these markets and treat $\hat{c}_{0,m}$ as an upper bound on the true hiding cost accordingly.

Under this convention, the concealment cost is economically material: 0.034 log points (approximately 3.4\% of property value) for apartments and 0.045 log points (4.5\%) for houses. That the two segments land within roughly one percentage point of each other suggests the cost is governed primarily by the legal disclosure regime and the state of registry coverage rather than by property-specific characteristics. The reduced-form premium provides an independent check: $\hat c_{0,m} \leq \hat\Delta_m$ holds market by market by construction, with equality exactly in the interior markets, so the inversion cannot exceed the bound the DiD imposes on it.

The contrast between property types is itself informative. Houses have an energy-efficiency premium three times larger than apartments ($\hat\beta_{1,m} \approx 0.051$ vs.\ 0.017 log points), reflecting their larger surface areas and older building stock, which make heating costs directly material to buyers' running expenses. A higher forgone premium makes concealment more expensive for low-quality sellers, and the data reflect this exactly: the hiding rate $\hat\sigma_{0,m}$ is well under half as large for houses (5.7\% vs.\ 14.0\%), and the share of strategic hiders is correspondingly lower (1.1\% vs.\ 3.5\% of all listings). The non-disclosure premium $\hat\Delta_m$ is also lower for houses (0.034 vs.\ 0.069), since a thinner pool of strategic hiders means buyers' beliefs about withheld quality are less pessimistic. The entire pattern (higher quality premium, lower hiding rate, lower non-disclosure premium) is the qualitative prediction of the model under equilibrium comparative statics across property types.

One further check concerns the gap between the non-disclosure premium $\hat\Delta_m$ and the structural hiding cost $\hat{c}_{0,m}$. In the apartment markets where both are available, the raw DiD estimate exceeds the inverted hiding cost by a factor of roughly two. This systematic gap is not a contradiction: it is what the model predicts. Uninformed sellers (approximately 83\% of non-disclosers for apartments and 89\% for houses) earn a non-disclosure premium simply by being in the market, pooling with strategic hiders and inflating the undisclosed category's average price above what strategic concealment alone would generate. $\hat\Delta_m$ is the market-level average of this pooled premium; $\hat{c}_{0,m}$ is the cost borne by the marginal informed hider, recovered by targeting only the informed-seller hiding rate $\hat\sigma_{0,m}$ in the equilibrium inversion. The two objects measure different things, and the fact that $\hat\Delta_m$ greatly exceeds $\hat{c}_{0,m}$ is precisely the compositional mechanism the model is designed to capture.
